\documentclass[11pt, a4paper]{article}

\usepackage{fontspec}
\usepackage{geometry}
\usepackage{fancyhdr}
\usepackage[hidelinks]{hyperref}
\usepackage[normalem]{ulem}
\usepackage{multicol}

\geometry{
  top=2cm,
  bottom=2cm,
  left=2cm,
  right=2cm,
  marginparsep=4pt,
  marginparwidth=1cm
}

\renewcommand{\headrulewidth}{0pt}
\pagestyle{fancyplain}
\fancyhf{}
\lfoot{}
\rfoot{}

\setlength{\parindent}{0pt}
\setlength{\parskip}{0pt}

\usepackage{xunicode}
\defaultfontfeatures{Mapping=tex-text}

\begin{document}

\begin{center}
\textbf{DSST289: Introduction to Data Science --- Fall 2026}
\end{center}

\vspace{0.5cm}

\textbf{Website}: \texttt{https://taylor-arnold.github.io/courses/dsst289-f26}

\bigskip

\textbf{Instructors:}
This course is co-taught by Prof.~Taylor Arnold (\texttt{taylor.arnold@richmond.edu}) and Prof.~Sam Kellogg (\texttt{sam.kellogg@richmond.edu}). Both of us will be in class and either of us is happy to hear from you about anything to do with the course material, the readings, the homework, or the project. Prof.~Arnold is the instructor of record, so please direct any questions about grades to him.

\bigskip

\textbf{Topics:}
Methods for collecting, manipulating, visualizing, exploring, and presenting
data.

\bigskip

\textbf{Format:}
The semester is broken into three units, each ending with an in-class exam, followed by a final project that runs through the end of the term. Class meetings are hands-on and interactive. Please bring a computer, pencil, and something to write on to each class. Class materials can be found on the course website.

\bigskip

\textbf{Software:}
We will be programming in the open-source Python programming language. Classwork can be run through Google's Colab environment in a web browser using a University email login. Each student will also have a personal shared Box folder, which we use only to return graded exams and to submit the final project. Other hardware and software may be used for some data collection tasks; all of these will be provided or freely accessible.

\bigskip

\textbf{Reading and Homework:}
Most class meetings will have a short reading and homework task posted on the course website. Homework is due at the start of class. We do not grade it for correctness, and it carries no points of its own, but making a full, genuine attempt at it is required to receive attendance credit for that class meeting. We will use it in class, and the exams and the final project are built directly on top of it. Coming to class without it means sitting out the part of the class that matters most.

\medskip

Please write your homework out by hand, on a fresh sheet of paper. Do not type it, print it, or fill in the margins of a printed handout. The exams are hand-written, and writing the code out yourself is most of the reason the homework helps. From time to time we will collect the homework at the start of class. When we do, we are only checking that you have done it and made a genuine attempt at it. 

\bigskip

\textbf{Attendance:}
At the beginning of each class meeting you will complete a brief class form that records your attendance and asks you to confirm that you have done that day's homework. \textit{Do not fill out the form if you are not present.} Counting as present requires both being in class and having made a full attempt at the homework, which the occasional collection described above lets us verify. Your attendance grade is the percentage of class meetings you attend, with your first two missed meetings forgiven: it is computed as the number of meetings you attended divided by the total number of meetings minus two, capped at 100. The two forgiven absences are there to cover routine illness and other unavoidable circumstances, so you do not need to explain them to us. If you find yourself needing to miss more than two meetings, contact us as soon as you can so that we can make appropriate arrangements.

\bigskip

\textbf{Exams:}
There are three in-class exams, given on the dates indicated on the course website. They consist of short, hand-written blocks of code, along with a small number of short written definitions on the third exam. A study guide for each exam, listing exactly what it covers and giving practice questions with answers, will be posted on the course website at least one week ahead of time. There is no final exam.

\bigskip

\textbf{External Speakers:}
You are required to attend at least two external talks over the course of the semester. These are a required component of the course but are not associated with a numerical grade. A list of talks is posted on the course website; any of those will count with not explaination needed. You may also count any other on-campus talk with a data science component. For a talk that is not on our list, you will be asked to write a short justification explaining what the data science component was. Attendance is recorded through the talk form linked on the course website.

\bigskip

\textbf{Final Project:}
The final project runs from late October through the end of the semester. You will
design a relational dataset about your own life, collect it by hand over a fixed
fourteen-day window, document it, check it, analyze it, and present it to the class. The full description, including the schedule and the deliverables, is posted on the course website. Two class meetings are set aside as project workshops, and the last three meetings of the semester are given over to presentations.

\bigskip

\textbf{Artificial Intelligence:}
You may use generative AI tools (ChatGPT, Claude, Copilot, and so on) for work done outside of class, as long as you document that you used them. Documentation means a short note attached to the work saying which tool you used and what you used it for. Nothing we do in class requires AI. The in-class work, the exams, and the hand-written homework are all meant to be done without it. If you find that you need a model to get through them, please come talk to us, because that usually means some piece of the material has not landed yet and we would rather hear about it early. You are also responsible for anything you submit under your name, including text or code that a model wrote for you.

\bigskip

\textbf{Getting Help:}
We will usually have time in class to answer questions about the course
material. If questions arise outside of class, please feel free to send
these by email to either of us. We are happy to schedule office hours for any extended questions or personal concerns. Because we know everyone has a busy schedule, we offer office hours primarily by appointment.

\bigskip

\textbf{Final Grades:}
The final grade for the semester is computed as follows:

\medskip

\begin{tabular}{ll}
Exam I & 20\% \\
Exam II & 20\% \\
Exam III & 20\% \\
Attendance & 20\% \\
Final Project & 20\% \\
\end{tabular}

\medskip

A letter grade is assigned as follows:
             A (93--100), A- (90--92),
B+ (87--89), B (83--86),  B- (80--82),
C+ (77--79), C (73--76),  C- (70--72), and F (0--69).

\end{document}
